\section{Comprensión del problema}

\subsection{El problema raíz}

Transportes Curimón S.A. asume el cien por ciento de la responsabilidad sobre la carga que
transporta, sobre las condiciones del viaje ante la autoridad y sobre el siniestro ante su
aseguradora. Opera esa responsabilidad con recursos que en su mayoría no controla. De los 374
tractocamiones que gestiona, 226 pertenecen a 148 transportistas subcontratados. De los 454
conductores que trabajan bajo su programación, 258 no son sus trabajadores
\parencite{bases_caso10_2026}.

Bajo el régimen de subcontratación de la Ley N.\textordmasculine{} 20.123 \parencite{ley20123}, esa
distancia no es un inconveniente administrativo. Es una exposición legal cuya magnitud la compañía
no ha cuantificado, y la advertencia del gerente general recogida en el acta de directorio fija el
límite de cualquier solución posible. El sesenta por ciento de la capacidad no pertenece a la
empresa y esas personas no son sus trabajadores, de modo que toda propuesta que suponga la
capacidad de darles una orden fracasa el primer día.

\subsection{Dimensión de la operación}

\begin{tablaAudit}{Volumetría operacional entregada por el mandante. Numeral 14.1 del Caso}{volumetria}{p{7.4cm}rr}
\textbf{Magnitud} & \textbf{Hoy} & \textbf{A tres años} \\ \midrule
Tractocamiones gestionados & 374 & 430 \\
Flota propia & 148 & 170 \\
Flota subcontratada & 226 & 260 \\
Transportistas subcontratados & 148 & 175 \\
Conductores bajo programación & 454 & 520 \\
Conductores externos & 258 & 300 \\
Semirremolques propios & 210 & 245 \\
Viajes al año & 96.000 & 118.000 \\
Kilómetros al año & 41.000.000 & 50.000.000 \\
Toneladas al año & 2.400.000 & 2.900.000 \\
Documentos electrónicos de transporte al año & 128.000 & 157.000 \\
Clientes activos & 84 & 100 \\
Puntos distintos de carga y descarga & 1.400 & 1.700 \\
Cruces fronterizos al año & 1.900 & 2.400 \\
Tractocamiones con telemetría de fábrica & 61 & 110 \\
Fechas de vencimiento vivas & 6.000 & 7.000 \\
\end{tablaAudit}

La red se despliega sobre un corredor de 3.000 kilómetros con cinco terminales, dos talleres
propios y un estanque de combustible. El 26 por ciento de los kilómetros se recorre en vacío.

\subsection{Fragilidad financiera}

La facturación anual alcanza los 78.000 millones de pesos con un margen operacional de 9 por
ciento. Tres de los ocho contratos principales se sirven bajo la línea de costo y representan en
conjunto el 31 por ciento del ingreso. El peor de ellos acumula cuatro años consecutivos a menos 14
por ciento, subsidiado por rutas rentables mediante un prorrateo ciego por ingreso. La
concentración comercial agrava el cuadro, porque ocho clientes de una cartera de 84 generan el 71
por ciento de la facturación.

Ese margen de 9 por ciento condiciona el diseño. La restricción 14 del Caso obliga a evaluar la
propuesta económica con especial atención al costo de operar durante 36 meses, de modo que toda
decisión con costo recurrente por unidad se multiplica por 374 y luego por 430.

\subsection{Los tres eventos de 2026}

\begin{enumerate}
    \item \textbf{14 de febrero.} Un camión subcontratado sale de la calzada en la Ruta 5 Sur. El
    conductor había manejado el día anterior para otra empresa y no había completado su descanso.
    Los registros del mandante mostraban once horas sin conducir para ella, y no tenían forma de
    mostrar qué ocurrió en esas once horas.
    \item \textbf{Abril.} Una fiscalización inmoviliza un tractocamión con sustancias peligrosas. La
    documentación no correspondía a la carga y el curso obligatorio del conductor estaba vencido
    hacía tres semanas. El vencimiento vivía en una planilla que se actualiza cuando alguien se
    acuerda, situación que el Decreto Supremo N.\textordmasculine{} 298 no admite \parencite{ds298}.
    \item \textbf{Junio.} El primer costeo por kilómetro y por ruta revela los subsidios cruzados
    entre contratos.
\end{enumerate}

Los tres comparten causa. La información que habría evitado cada uno existía, repartida entre tres
plataformas de posicionamiento, una telemetría que nadie descarga, una liquidación de combustible
que llega con 40 días de atraso y papeles que viajan en la cabina.

\subsection{El estado de los datos}

\begin{tablaAudit}{Brechas de información verificadas en el Capítulo 7 del Caso}{brechas}{p{8.6cm}p{5.4cm}}
\textbf{Indicador} & \textbf{Situación actual} \\ \midrule
Conductores con control de jornada & 196 de 454 \\
Jornada previa de un conductor de tercero & Inexistente \\
Tacógrafos digitales cuya información se descarga & Cero \\
Tractocamiones con telemetría de fábrica sin descargar & 61 \\
Camiones sin dispositivo de posicionamiento & 34 \\
Plataformas de posicionamiento incompatibles & 3, una sin exportación \\
Fechas de vencimiento vivas en planillas separadas & 6.000 en 4 planillas \\
Espera media en punto de carga & 3 horas 10 minutos \\
Cobros por espera objetados por el cliente & 71 por ciento \\
Documentos de respaldo que llegan incompletos o no llegan & 4,2 por ciento \\
Liquidación mensual a 148 transportistas & 9 días, 8 personas, 11 por ciento se corrige \\
Detenciones por sobrepeso en 2025 & 142 \\
\end{tablaAudit}

\subsection{La exigencia de 2029}

El cliente exportador mayor, responsable del 19 por ciento de los ingresos, condicionó la
renovación de su contrato a cuatro exigencias. Documento electrónico de transporte integrado punta
a punta. Posición de la carga en tiempo real. Emisiones por tonelada kilómetro verificadas por un
tercero conforme a ISO 14083 \parencite{iso14083} y al marco GLEC \parencite{glec2023}. Y acreditación del cumplimiento de
jornada en cada viaje, incluidos los camiones subcontratados.

La cuarta exigencia es el problema central del Caso. El régimen del Artículo 25 bis del Código del
Trabajo \parencite{codigo_trabajo_25bis} recae sobre el empleador, y para los 258 conductores
externos el empleador es el transportista, no el mandante. Esa constatación no es un obstáculo. Es
el anclaje jurídico de la solución que la Sección 3 desarrolla, porque confirma que la obligación
se persigue por el contrato con el dueño del camión y no con el conductor.

\subsection{Los nodos donde ocurre la operación}

\begin{tablaAudit}{Nodos operacionales y su condición. Capítulo 3 del Caso}{nodos}{p{4.2cm}p{9.8cm}}
\textbf{Nodo} & \textbf{Condición determinante para el diseño} \\ \midrule
Terminal San Bernardo & Casa matriz, torre de programación 24x7, taller principal. Único punto donde converge todo y donde se instala cualquier equipamiento a bordo \\
Cuatro terminales regionales & Antofagasta, Talca, Los Ángeles y Puerto Montt. Enlace de un proveedor, sin respaldo en tres de los cuatro \\
La cabina del camión & Lugar de trabajo real. Ninguna interacción admisible durante la marcha \\
La ruta & Tramos de más de 80 kilómetros continuos sin cobertura. Ocurre allí el cien por ciento del riesgo y existe allí la menor visibilidad \\
Puntos de carga y descarga & Instalaciones de terceros con reglas y sistemas propios. No se puede instalar equipamiento \\
Plazas de pesaje & El cliente carga y a la compañía la pesan. Un sobrepeso es multa e inmovilización \\
Paso Los Libertadores & Cierra por nieve entre junio y septiembre, en episodios de hasta 12 días continuos \\
Talleres en ruta & Proveedores externos cuya intervención hoy no queda registrada \\
\end{tablaAudit}
